% ============================================================================
%  系统开发工具基础 · 实验报告模板
% ----------------------------------------------------------------------------
%  设计依据：课程《系统开发工具基础》（周小伟，计算机科学与技术学院）
%  实验报告须包含：
%     (1) 课上给助教检查的内容和结果
%     (2) 课后练习内容和结果（每次课 >10 个实例）
%     (3) 解题感悟等
%  另须将 GitHub/Gitee 仓库链接与 commit 记录截图写入实验报告。
%
%  编译方式（推荐，XeLaTeX + shell-escape 以启用 minted 代码高亮）：
%     latexmk -xelatex -shell-escape -halt-on-error 实验报告模板.tex
%  或单次编译：
%     xelatex -shell-escape -interaction=nonstopmode 实验报告模板.tex
%
%  使用说明：
%     1. 在下方「基本信息」处填写姓名、学号、周次、日期。
%     2. 在「一、课上实验检查」中，为每道题新建一个 experiment 环境。
%     3. 在「二、课后练习」中，用 exercise 环境逐条记录练习内容与结果。
%     4. 在「三、解题感悟」中总结本讲收获与思考。
%     5. 在「四、版本控制与提交记录」中填写仓库链接并插入 commit 截图。
% ============================================================================

\documentclass[UTF8, zihao=-4]{ctexart}

% ------------------------------- 基本信息 ---------------------------------
% 每次写报告只需修改这里
\newcommand{\university}{中国海洋大学}
\newcommand{\coursename}{系统开发工具基础}
\newcommand{\teacher}{周小伟}
\newcommand{\studentname}{姓名}
\newcommand{\studentid}{学号}
\newcommand{\reportweek}{第~1~周}
\newcommand{\reportdate}{2026 年 8 月 24 日}

% ------------------------------- 宏包加载 ---------------------------------
\usepackage[a4paper, top=2.5cm, bottom=2.5cm, left=2.8cm, right=2.8cm]{geometry}
\usepackage{amsmath, amssymb}
\usepackage{graphicx}
\usepackage{booktabs}
\usepackage{array}
\usepackage{xcolor}
\usepackage{float}                 % 提供 [H] 图片排版
\usepackage{caption}
\usepackage{fancyhdr}
\usepackage[colorlinks=true, linkcolor=cprimary, urlcolor=csecondary, citecolor=csecondary]{hyperref}
\usepackage[most]{tcolorbox}
\usepackage{minted}                % 代码高亮（需 -shell-escape）
\usepackage{enumitem}

% ------------------------------- 颜色定义 ---------------------------------
\definecolor{cprimary}{RGB}{15, 76, 129}       % 主色：深蓝
\definecolor{csecondary}{RGB}{46, 134, 193}    % 辅色：亮蓝
\definecolor{caccent}{RGB}{230, 126, 34}       % 强调：橙
\definecolor{clight}{RGB}{240, 246, 251}       % 浅蓝底
\definecolor{codebg}{RGB}{248, 250, 252}       % 代码块底色
\definecolor{cgray}{RGB}{90, 90, 90}           % 灰色文字

% ------------------------------- 字体（可选）-------------------------------
% 如需自定义中文字体可取消注释，例如：
% \setCJKmainfont[AutoFakeBold=true]{SimSun}
% \setCJKsansfont{Microsoft YaHei}

% ------------------------------- 代码块样式 -------------------------------
\usemintedstyle{vs}
\setminted{
  breaklines=true,
  fontsize=\small,
  bgcolor=codebg,
  frame=single,
  rulecolor=csecondary!35,
  framesep=4pt,
  autogobble=true,
  baselinestretch=1.1,
}

% ------------------------------- 页眉页脚 ---------------------------------
\pagestyle{fancy}
\fancyhf{}
\fancyhead[L]{\small\textcolor{cprimary}{\coursename}}
\fancyhead[R]{\small\textcolor{cprimary}{\reportweek \,实验报告}}
\fancyfoot[C]{\small\thepage}
\renewcommand{\headrulewidth}{0.6pt}
\renewcommand{\headrule}{\hbox to\headwidth{\color{cprimary}\leaders\hrule height \headrulewidth\hfill}}
\renewcommand{\footrulewidth}{0pt}

% ------------------------------- 自定义命令 -------------------------------
% 报告封面标题
\newcommand{\makereporttitle}{%
  \begin{center}
    {\small\color{cgray}\university\par}
    \vspace{0.8em}
    {\Huge\bfseries\color{cprimary}\coursename\par}
    \vspace{0.3em}
    {\Large\bfseries 实\quad 验\quad 报\quad 告\par}
    \vspace{0.6em}
    {\large\reportweek\par}
    \vspace{0.8em}
    {\color{cprimary}\rule{0.7\textwidth}{1.2pt}\par}
    \vspace{1.2em}
    \renewcommand{\arraystretch}{1.6}
    \begin{tabular}{>{\bfseries}r@{\hspace{1.5em}}l}
      姓\quad 名 & \studentname \\
      学\quad 号 & \studentid \\
      授课教师 & \teacher \\
      实验日期 & \reportdate \\
    \end{tabular}
  \end{center}
  \vspace{1em}
}

% 内容小标题（实验内容 / 关键命令 / 实验结果 等）
\newcommand{\contenthead}[1]{%
  \medskip\noindent\textcolor{cprimary}{\bfseries #1}\par\smallskip
}

% 截图命令：\screenshot[宽度]{图片路径}{图注}
\newcommand{\screenshot}[3][0.85\textwidth]{%
  \begin{figure}[H]
    \centering
    \includegraphics[width=#1]{#2}
    \caption{#3}
  \end{figure}
}

% GitHub/Gitee 仓库链接
\newcommand{\gitlink}[1]{\url{#1}}

% 空栏占位（提示待填写）
\newcommand{\todo}[1]{\textcolor{caccent}{\emph{（待填写：#1）}}}

% ------------------------------- 实验环境 ---------------------------------
% 课上实验：\begin{experiment}{题号}{主题} ... \end{experiment}
\newtcolorbox{experiment}[2]{%
  enhanced,
  breakable,
  colback=white,
  colframe=cprimary,
  colbacktitle=cprimary,
  coltitle=white,
  fonttitle=\bfseries,
  title={课上实验\quad#1\quad——\quad#2},
  attach title to upper={\par},
  before title={\smallskip},
  after title={\smallskip},
  left=4mm, right=4mm, top=3mm, bottom=3mm,
}

% 课后练习：\begin{exercise}{练习编号/名称} ... \end{exercise}
\newtcolorbox{exercise}[1]{%
  enhanced,
  breakable,
  colback=clight,
  colframe=csecondary,
  colbacktitle=csecondary,
  coltitle=white,
  fonttitle=\bfseries,
  title={课后练习\quad#1},
  attach title to upper={\par},
  before title={\smallskip},
  after title={\smallskip},
  left=4mm, right=4mm, top=3mm, bottom=3mm,
}

% ============================================================================
%                                正文
% ============================================================================
\begin{document}

\makereporttitle

% ---------------------------- 一、课上实验检查 ------------------------------
\section*{一、课上实验检查}

本部分记录每次课上完成并给助教检查的四道题目，包括\textbf{实验内容}、
\textbf{关键命令/代码}与\textbf{实验结果}（附截图）。

\begin{experiment}{第 1 题}{含空格文件名的批量整理（Shell 基础与文件系统）}
  \contenthead{实验内容}
  在 \texttt{q01} 中完全使用命令行创建含空格与隐藏文件的目录结构，长格式列出
  \texttt{input} 下全部项目，并将所有 \texttt{.txt} 文件复制到
  \texttt{work/<学号>/}，随后设置权限并生成清单。

  \contenthead{关键命令}
  \begin{minted}{bash}
    mkdir -p input/docs input/tmp
    printf 'alpha\nbeta\n' > 'input/docs/notes one.txt'
    printf 'hidden\n' > input/docs/.secret.txt
    touch input/tmp/empty.txt
    printf 'log line 1\nlog line 2\n' > input/run.log
    ls -la input
    # 复制并保留相对结构（注意文件名中的空格）
    cp --parents input/**/*.txt work/<学号>/ 2>/dev/null || find input -name '*.txt' -exec ...
  \end{minted}

  \contenthead{实验结果}
  \todo{粘贴终端运行结果截图与 \texttt{inventory.txt} 内容}
\end{experiment}

\begin{experiment}{第 2 题}{随机访问日志统计（Shell 管道与文本处理）}
  \contenthead{实验内容}
  编写 \texttt{analyze.sh} 统计 HTTP 5xx 次数最多的前 2 个 \texttt{path} 及平均
  \texttt{latency\_ms}。

  \contenthead{关键命令}
  \begin{minted}{bash}
    #!/usr/bin/env bash
    [ -f "$1" ] || { echo "error: no such file: $1" >&2; exit 1; }
    tail -n +2 "$1" | awk -F, '$4 ~ /^5[0-9][0-9]$/ {c[$3]++} END {for (p in c) print c[p], p}' \
      | sort -k1,1nr -k2,2 | head -2
    tail -n +2 "$1" | awk -F, '{s+=$5; n++} END {printf "%.2f\n", s/n}'
  \end{minted}

  \contenthead{实验结果}
  \todo{分别对 access.csv 与不存在的文件运行脚本，粘贴输出}
\end{experiment}

\begin{experiment}{第 3 题}{制造、解决并解释一次合并冲突（Version Control and Git）}
  \contenthead{实验内容}
  从零创建 Git 仓库，在 \texttt{main} 与两个特性分支上修改 \texttt{config.txt}，
  合并时解决冲突并保留指定内容。

  \contenthead{关键命令}
  \begin{minted}{bash}
    git init
    echo 'mode=normal' > config.txt && git add . && git commit -m 'init'
    git switch -c feature-a && echo 'mode=safe' > config.txt && git commit -am 'a'
    git switch main && git switch -c feature-b && echo 'mode=fast' > config.txt && git commit -am 'b'
    git switch main && git merge feature-a && git merge feature-b
    # 解决冲突：保留 mode=safe 并追加 note=reviewed
    git log --all --graph --decorate --oneline
  \end{minted}

  \contenthead{实验结果}
  \todo{粘贴冲突标记、解决后的 config.txt 与 git log 图}
\end{experiment}

\begin{experiment}{第 4 题}{修复并构建一页技术说明（LaTeX 文档编辑）}
  \contenthead{实验内容}
  补全 \texttt{report.tex}：加入带编号与 label 的公式、2×2 表格，并在正文中交叉引用。

  \contenthead{关键命令}
  \begin{minted}{latex}
    \begin{equation}\label{eq:sum}
      S = \sum_{i=1}^{n} i = \frac{n(n+1)}{2}
    \end{equation}
    由式~\eqref{eq:sum} 可知…… 结果见表~\ref{tab:data}。
  \end{minted}

  \contenthead{实验结果}
  \todo{粘贴生成的 report.pdf 截图，确认交叉引用无 “??”}
\end{experiment}

% ------------------------------ 二、课后练习 -------------------------------
\section*{二、课后练习}

每次课后完成 \textbf{10 个以上}实例，逐条记录\textbf{练习内容}与\textbf{结果}。

\begin{exercise}{1 · Git 基础命令}
  \contenthead{练习内容}
  练习 \texttt{git init / add / commit / log / diff / status} 等命令。

  \contenthead{结果}
  \todo{命令输出或小结}
\end{exercise}

\begin{exercise}{2 · Shell 常用命令}
  \contenthead{练习内容}
  练习 \texttt{ls / find / grep / awk / sed / chmod} 等命令。

  \contenthead{结果}
  \todo{命令输出或小结}
\end{exercise}

% 其余实例可复制上面的 exercise 环境继续添加（编号 3、4、…）

% ------------------------------ 三、解题感悟 -------------------------------
\section*{三、解题感悟}

\contenthead{本讲收获}
\todo{总结本讲学到的知识点，例如：管道与重定向、Git 分支与冲突解决、
LaTeX 交叉引用等。}

\contenthead{遇到的问题与解决}
\todo{记录遇到的困难、排查过程与最终解决方式。}

\contenthead{对课程的建议}
\todo{可选。}

% -------------------------- 四、版本控制与提交记录 --------------------------
\section*{四、版本控制与提交记录}

\contenthead{仓库链接}
GitHub/Gitee 仓库地址：\gitlink{https://github.com/yourname/your-repo}
（将 \texttt{yourname/your-repo} 替换为自己的用户名与仓库名）

\contenthead{提交记录截图}
\todo{插入 GitHub/Gitee 的 commit 记录截图，体现多次提交、非单次全量提交}
\begin{center}
  \fbox{\parbox{0.7\textwidth}{\centering\color{cgray}【在此插入 commit 记录截图】\par}}
\end{center}

\end{document}
